\section{Introduction}

Documents are one of the fundamental forms of information preservation and exchange. In each year, tens of millions of documents are created, read, saved, and dispatched~\cite{Bornmann2014GrowthRO}. Beyond unstructured pure-text, documents feature both complicated layout structures and information across distinct modalities such as text, table, chart, image, \etc Accordingly, the automatic understanding of documents (Document Understanding; DU) stands as a long-standing task in urgent and practical needs.

Recently, a number of LVLMs, both closed-source ones (GPT-4o~\cite{gpt4o}, Gemini-1.5~\cite{geminiteam2024gemini}, Claude-3~\cite{claude3}, \etc) and open-source ones (InternLM-XC2-4KHD~\cite{dong2024internlm}, InternVL-Chat~\cite{chen2024far}, Otter~\cite{li2023otter}, LLaVA-NeXT~\cite{li2024llavanext-strong}, CogVLM~\cite{wang2023cogvlm}, mPLUG-DocOwl 1.5~\cite{hu2024mplugdocowl}, TextMonkey~\cite{liu2024textmonkey}, \etc) have been developed and presented the great potential to handle documents. Most of them have achieved promising performance on single-page DU datasets like DocVQA~\cite{Mathew2020DocVQAAD}, ChartQA~\cite{masry-etal-2022-chartqa}, InfoVQA~\cite{Mathew2021InfographicVQA}, TAT-DQA~\cite{zhu2022towards}, \etc However, considerable amounts of documents in the real world are long-context documents with tens or even hundreds of pages. The understanding of these lengthy documents brings new challenges for LVLMs from at least two aspects: (1) \textbf{Localization}: identify and retrieve information from massive, heterogeneous information (similar to the \textit{needle in a haystack} task); (2) \textbf{Cross-page comprehension}: collect and reason over multi-source information across different pages. These two kinds of abilities are beyond the evaluation scopes of the aforementioned single-page DU datasets. Some recent DU datasets~\cite{tito2023hierarchical, VanLandeghem2023DocumentUD, SlideVQA2023} feature multiple-page DU, but almost all their documents are either as short of only several pages or of low information density, making the localization-related questions over-simple. Additionally, few (if any) questions in these datasets necessitate cross-page comprehension. See more detailed related work in Section~\ref{sec: related_work}. In summary, there lacks a unified and high-quality benchmark on lengthy documents, leaving the evaluation of long-context DU largely unexplored.

In this paper, we present \benchmarkname, a benchmark designed to evaluate the \textbf{M}ulti-\textbf{M}odality \textbf{Long}-context \textbf{Doc}ument understanding abilities of LVLMs. Towards a comprehensive benchmark, it incorporates lengthy documents from both four existing datasets~\cite{masry-etal-2022-chartqa, VanLandeghem2023DocumentUD, SlideVQA2023, islam2023financebench} and other various papers, brochures, \etc Consequently, our benchmark includes 135 PDF-formatted documents spanning across 7 diverse domains, with each document averaging 47.5 pages and 21,214.1 textual tokens. Regarding the questions, we employ ten expert-level annotators to (1) edit questions associated with documents from existing datasets to meet our benchmark's standard and (2) create new questions for all collected documents to expand the scale of the benchmark. Then a three-round, semi-automatic reviewing process ensures the benchmark's annotation quality. As a result, \benchmarkname comprises 1,082 human-annotated questions, with 184 sourced from four existing datasets and 898 newly annotated. Being a multi-modal benchmark, the answer to each question requires evidence from one or more of these five in-document sources: \textit{text}, \textit{layout}, \textit{chart}, \textit{table}, and \textit{image}. Questions are categorized into three types based on the number of evidence pages~\footnote{Given a document $D$ and a question $q$ upon $D$, We call page $P$ (in document $D$) an \textit{evidence page} of $q$ if the answer of $q$ necessitates one or more pieces of evidence in page $P$.}, with examples illustrated in Figure~\ref{fig: examples}(a): (1) 494 \textit{single-page} questions (with one evidence page) mainly to evaluate localization abilities, (2) 365 \textit{cross-page} questions (with multiple evidence pages) to assess cross-page comprehension, and (3) 223 \textit{unanswerable} questions (no evidence for answering it, \ie no evidence pages) to reduce shortcuts and measure LVLMs' potential hallucinations. Meta-information including evidence pages, sources, and answer formats, is preserved for fine-grained evaluation and analysis. Detailed descriptions of the annotation pipeline and statistics can be found in Section~\ref{sec: dataset}.

We conduct extensive experiments on \benchmarkname to evaluate the long-context DU abilities of 14 LVLMs, including 4 proprietary and 10 open-source ones. Given a document, we screenshot each page and feed all of these PNG-formatted images to LVLMs in an end-to-end approach. For comparison, we also convert the documents to textual format by optical character recognition (OCR) and evaluate another 6 proprietary and 4 open-source 10 LLMs (6 proprietary and 4 open-source ones). The results in Figure~\ref{fig: examples}(c) highlight the challenges that current LVLMs face with long-context DU. The best-performing LVLM, GPT-4o, achieves an overall F1 score of only 44.9\%, while the second-best LVLM, GPT-4V, scores 30.5\%. Moreover, all the remaining LVLMs tested with multi-modal documents performed worse than single-modal LLMs handling lossy, OCR-parsed texts. Specifically, the Gemini-1.5-Pro and Claude-3-Opus present 4.2\% and 6.4\% absolute decrease when the inputs change from document screenshots to OCR-parsed texts. Regarding open-source models, the best-performing LVLM lags behind the best-performing LLM by 11.7\%. These results reveal that long-context DU is a far-from-resolved task for current LVLMs.

\begin{table}[t]
\centering
\small
\caption{Comparison between our benchmark and previous DU datasets. \textbf{Unans.}: unanswerable question. \textbf{TXT/L/C/TAB/I}: pure text/generalized layout/chart/table/image. \textbf{Doc. Rel.}: document relevance. Whether document information is indispensable for the answer. \textbf{Avg. Position}: the average page index on which the answer evidence is located. *:Statistics from \cite{Borchmann2021DUEED}.}
\resizebox{0.95\linewidth}{!}{
\begin{tabular}{l|cc|ccccc}
\toprule
\multirow{2}{*}{\textbf{Benchmarks}}  & \multicolumn{2}{c}{\textbf{Document}} & \multicolumn{2}{|c}{\textbf{Question type}} & \multicolumn{3}{c}{\textbf{Answer Evidence}} \\
& \# Pages & \# Tokens & Cross-page (\%) & Unans. (\%) & Doc. Rel. & Source & Avg. Position  \\
\midrule
DocVQA~\cite{Mathew2020DocVQAAD} & 1.0 & 151.5 &  \xmark & \xmark & \cxmark &  TXT/L/C/TAB/I & - \\
ChartQA~\cite{masry-etal-2022-chartqa}\tablefootnote{We view website screenshots and posters as generalized documents and define \textit{equivalent page number} (\texttt{EPN}) to measure their context lengths: $\texttt{EPN(D) = ceil}(\frac{\texttt{Pixel(D)}}{P})$. Here \texttt{Pixel(D)} is the pixel number of generalized document \texttt{D}, and \texttt{P} is the average pixel numbers of each page (converting from .pdf to .png format with resolution 240) in \benchmarkname. \label{equivalent pages}} & 1.0 & 236.9 & \xmark & \xmark & \cmark & C & -  \\
InfoVQA~\cite{Mathew2021InfographicVQA}\textsuperscript{\ref{equivalent pages}} & 1.2 & 288.0  & \xmark & \xmark & \cxmark &  L/C/TAB/I & - \\
TAT-DQA~\cite{zhu2022towards} & 1.1 & 577.0 & \xmark & \xmark & \cxmark & TXT/TAB & -   \\
VisualWebBench~\cite{liu2024visualwebbench}\textsuperscript{\ref{equivalent pages}} & 1.0 & 452.4  & \xmark & \xmark & \cmark & LAY/I & - \\
PWC~\cite{kardas-etal-2020-axcell} & \textasciitilde 12* & \textasciitilde 7000* & \xmark & \xmark & \cxmark & TAB & - \\
MP-DocVQA~\cite{tito2023hierarchical}  & 8.3 & 2026.6 & \xmark & \xmark & \cxmark & TXT/L/C/TAB/I & 6.0 \\
DUDE~\cite{VanLandeghem2023DocumentUD} & 5.7 & 1831.5 & \cmark (2.1\%) & \cmark(12.7\%) & \cxmark & TXT/L/C/TAB/I & 2.5 \\
SlideVQA~\cite{SlideVQA2023} & 20.0 & 2030.5 & \cmark (13.9\%) & \xmark & \cxmark & TXT/L/C/TAB/I & 9.1 \\
\midrule
\benchmarkname & 47.5 & 21214.1 & \cmark (33.0\%) & \cmark (22.5\%) & \cmark & TXT/L/C/TAB/I & 23.6 \\
\bottomrule
\end{tabular}}
\label{tab:dataset_comparison}
\end{table}